% Source: https://artofproblemsolving.com/community/c3h1281255_close_to_mathcounts_2012_team_10
% Stable identity: aops-topic-1281255-q01
Take points on each line at the same distance from the origin: unit vectors along the lines are $\left(\dfrac{5}{13},\dfrac{12}{13}\right)$ and $\left(\dfrac{3}{5},\dfrac{4}{5}\right)$ (from the $5$--$12$--$13$ and $3$--$4$--$5$ triangles). The sum of two equal-length vectors is the diagonal of a rhombus, which bisects the angle between them.

\begin{asy}
usepackage("amsmath");
size(4.5cm);
defaultpen(fontsize(10pt));
draw((-0.3,0)--(5,0),gray,Arrow); draw((0,-0.3)--(0,6.4),gray,Arrow);
pair U=4*(5/13,12/13), V=4*(3/5,4/5); draw((0,0)--1.3*U,red); draw((0,0)--1.5*V,blue); draw((0,0)--(3.5,6.125),purple+dashed);
dot(U); dot(V); draw(U--U+V--V,gray+dotted); dot(U+V);
label("$y=\dfrac{12}{5}x$",1.3*U,W,red); label("$y=\dfrac{4}{3}x$",1.5*V,E,blue);
label("\mbox{}",(min(currentpicture,user=true)+max(currentpicture,user=true))/2,black); // keeps the text after the figure black when LaTeX typesets the labels inline
\end{asy}

The sum is
\[
\left(\dfrac{5}{13}+\dfrac{3}{5},\ \dfrac{12}{13}+\dfrac{4}{5}\right)=\left(\dfrac{64}{65},\dfrac{112}{65}\right),
\]
so the bisector has slope $\dfrac{112}{64}=\boxed{\dfrac{7}{4}}$. (Both direction vectors point into the first quadrant and make an acute angle, so this is the acute-angle bisector.)
